\documentclass[10pt]{article}

\usepackage[utf8]{inputenc}

\usepackage[T1]{fontenc}

\usepackage{amsmath}

\usepackage{amsfonts}

\usepackage{amssymb}

\usepackage[version=4]{mhchem}

\usepackage{stmaryrd}

\usepackage{graphicx}

\usepackage[export]{adjustbox}

\graphicspath{ {./images/} }

\usepackage{bbold}



\begin{document}

Теорема аппроксимации может служить отправным пунктом определения различных классов П. П. Ф. При әтом аппроксимирующие полиномы $P(x)$ могут содержать "посторонние», т. е. отличные от показателей Фурье функции $f(x)$, показатели. Однако для нек-рых приложений теоремы аппроксимации важен тот факт, что можно совершенно избежать введения в $P(x)$ показателей, отличных от показателей Фурье функции $f(x)$.



В связи с представимостью П. П. ф. обобщенными рядами Фурье возникает вопрос о признаках сходимости для этих рядов и приобретают значение разнообразные методы суммирования обобщенных рядов Фурье (метод Бохнера - Фейера и др.). Так, получены признак абсолютной сходимости обобщенных рядов Фурье с линейно независимыми показателями Фурье, признак равномерной сходимости рядов Фурье, когда $\left|\lambda_{k}\right| \rightarrow \infty$ при $k \rightarrow \infty$, и аналогичный признак в случае, когда $\lambda_{k} \rightarrow 0$ при $k \rightarrow \infty$.



Значение признаков равномерной сходимости в теории П. П. ф. подчеркивается следующей теоремой: если тригонометрич. ряд $\sum_{k} a_{k} e^{i \lambda_{k}^{x}}$ сходится равномерно на всей действительной оси, то он является рядом Фурье своей суммы $S(x) \subset U$-П. П. Следствие: существуют равномерные П. П. ф. с произвольным счетным множеством показателей Фурье. В частности, показатели Фурье равномерной П. П. Ф. могут иметь предельные точки на конечном расстоянии или даже располагаться всюду плотно.



Кроме понятия замыкания или почти периода, для өпределения П. п. ф. можно использовать понятие сдвига. Так, функция $f(x)$ будет равномерной П. п. ф. тогда и только тогда, когда из каждой бесконечной последовательности функций $f\left(x+h_{1}\right), f\left(x+h_{2}\right), \ldots$, где сдвиги $h_{1}, h_{2}, \ldots$ - произвольные действительные числа, можно выбрать равномерно сходящуюся подпоследовательность. Это определение служит отправной точкой при рассмотрении П. П. ф. на группах.



Основные факты теории П. П. Ф. остаются справедливыми и в том случае, если рассматривать понятие обобщенного сдвига. Возможны и полезны другие обобщения. П. П. Ф. со значениями в евклидовом $n$-мерном пространстве, в банаховом или метрич. пространстве, аналитические и гармонические П. П. Ф.



Лит.: [1] Б ор Г., Почти периодические функции, пер. c Hem., M.-J., 1934; [2] B e s i c o v i t c h A. S., Almost



\includegraphics[max width=\textwidth, center]{2023_07_08_a8574ed425b8d913ed3eg-1(1)}

периодатем. наук", 1968 , т. 23 , в. 4, c. $117-78$; [5] R u d i n "W. пеxи мaтeм. науk", 1968, T. 23, в. 4, c. $117-78 ;[5] \mathrm{R}$ u d i n W.,

Fourier analysis on groups, N.Y.-L., $1962 ;[6]$ J е в и T a H Б М., Операторы обобщенного сдвига и некоторые их применения, М., $1962 ;$ [7] К р а с н о с е л b с к и й М. А., Б у p д колебания, М., 1970 . E. А. Бредихина.



ПОчТИ ПЕРИОДИЧЕСКАЯ ФУНКЦИЯ АНАЛИТИЧЕСКАЯ - аналитическая функция $f(s), s=\sigma+i \tau$, регулярная в полосе $-\infty \leqslant \alpha<\sigma<\beta \leqslant+\infty$ и разложимая в ряд вида



\[

\sum a_{n} e^{i \lambda_{n} s}

\]



где $a_{n}$ - комплексные, $\lambda_{n}$ - действительные числа. Действительное число $\tau$ наз. $\varepsilon-\Pi$ о т т и п е р и о д о м функции $f(s)$, если во всех точках полосы $(\alpha, \beta)$ выполняется неравенство



\[

|f(s+i \tau)-f(s)|<\varepsilon .

\]



П. п. Ф. а.- аналитическая функция, регулярная в полосе $(\alpha, \beta)$ и обладающая для каждого $\varepsilon>0$ относи-



\includegraphics[max width=\textwidth, center]{2023_07_08_a8574ed425b8d913ed3eg-1}

гично определяется П. п. ф. а. в замкнутой полосе $\alpha \leqslant \sigma \leqslant \beta$. П. П. ф. а. в полосе $[\alpha, \beta]$ на каждой прямой полосы является равномерной почти периодич. функцией от действительного переменного $\tau$, она ограничена в $[\alpha, \beta]$, т. е. в любой внутренней полосе. Если функция $f(s)$, регулярная в полосе $(\alpha, \beta)$, является равномерной почти периодич. функцией хотя бы на одной единственной прямой $\sigma=\sigma_{0}$ этой полосы, то ограниченность $f(s)$ в $[\alpha, \beta]$ влечет за собой ее почти периодичность во всей полосе $[\alpha, \beta]$. В результате теория П. п. ф. а. оказывается в основе своей аналогичной теории почти периодических функций от действительного переменного. Поэтому на П. П. Ф. а. легко переносятся многие важные факты последней теории: теорема единственности, равенство Парсеваля, правила действий над рядами Дирихле, аппроксимационная теорема и ряд др. теорем.



Лит.: [1] Б о p $\Gamma$, Почти перводические функции, пер. с нем., М.- Л., 1934; [2] Л е в и т а н Б. М., Почти-периоди-

ческие функции, М., 1953.



ПОЧТИ ІЕРИОДИЧЕСКАЯ ФУНКЦИЯ НА ГРУІПЕ - обобщение почти периодич. функций, заданных на $\mathbb{R}^{1}$.



Пусть $G$ - (абстрактная) группа. Комплекснозначная функция $f(x), x \in G$, наз. П р а в о й по ч т и п ери и ди ч е ск о й функци е й (п. п. ф.), если семейство $f(x a)$, где $a$ пробегает всю группу $G$, компактно в смысле равномерной сходимости на $G$, т. е. из каждой последовательности $f\left(x a_{1}\right), f\left(x a_{2}\right), \ldots$ можно выделить равномерно на $G$ сходящуюся подпоследовательность. Аналогично определяются ле в а я п о ч т и пе р и о ди ч е с к я фу у кция на группе $G$. Оказывается, что всякая правая (левая) п. п. ф. является одновременно левой (правой) п. П. ф., и имеет место компактность семейства $f(a x b)$, где $a, b$ независимо пробегают группу $G$. Последнее свойство часто принимается в качестве определения п. п. ф. на $G$. Совокупность всех П. П. Ф. на $G$ есть банахово пространство с нормой $\|f\|=\sup _{x \in G}|f(x)|$.



Теория п. п. ф. на группе существенно опирается на те о ре му о с р е днем значен и и (см. [5], [8]). Линейный функционал $M_{x}\{f(x)\}$, заданный на пространстве всех п. п. ф., наз. с р е дн и м 3 на ч ен и е м, если



\begin{enumerate}

  \item $M_{x}\{1\}=1, M_{x}\{f(x)\} \geqslant 0$ для $f(x) \geqslant 0$ и $M_{x}\{f(x)\}>0$ для $f \not 00$; $M_{x}\{f(x a)\}=M_{x}\{f(a x)\}=M_{x}\left\{f\left(x^{-1}\right)\right\}$.

\end{enumerate}



Унитарная матрица $g(x)=\left\{g_{i j}(x)\right\}_{i, j=1}^{r}$, заданная на $G$, наз. у ни т а р ны м пре д с т а в лен и е м гр у п п ы $G$, если $g(e)=I_{r}(e-$ единица группы $G$, $I_{r}$ - единичная матрица порядка $r$ ) и для любых элементов $x, y \in G$ имеет место равенство $g(x y)=g(x) g(y)$. Число $r$ наз. р а зме р ност ь ю п р е д с т а в л ен и я $g$. Матричные әлементы $g_{i j}(x)$ суть п. п. ф. на G. В теории п. п. Ф. на группе они играют ту же роль, что и функции $\exp (i \lambda x)$ в теории п. п. ф. на прямой $\mathbb{R}^{1}$,



Два представления $g(x)$ и $g^{\prime}(x)$ наз. э к в и в а л е н тн ы и, если существует такая постоянная матрица $A$, что $g^{\prime}(x)=A^{-1} g(x) A$. Представление $g$ наз. н е п р ив о д и мым, если семейство матриц $g(x), x \in G$, не имеет общего нетривиального инвариантного подпространства в $\mathbb{R}$. Множество всех неприводимых унитарных представлений разбивается на классы эквивалентных между собой представлөний. Пусть из каждого класса эквивалентных представлений выбрано по одному представлению и полученное множество обозначено $S$. Тогда множество п. П. ф. на $G$



\[

H=\left\{\varphi_{\lambda}(x)\right\}=\left\{\varphi, \varphi=g_{i j}, g \in S\right\}

\]



оказывается ортогональной (хотя, вообце говоря, несчетной) системой.



Т о о е м а 1 (ра в енство П а р с е в а ля). Если для п. п. ф. $f(x)$ положить



\[

f \sim \sum \varphi_{\lambda} \frac{M_{x}\left\{f(x) \bar{\varphi}_{\lambda}(x)\right\}}{M_{x}^{1 / 2}\left\{\left|\varphi_{\lambda}^{(x)}\right|^{2}\right\}},

\]





\end{document}